\section{Satellite galaxies around a massive central}

\subsection{Computing $A$}
%# Write a numerical integrator based on some form of Richardson extrapolation to 
%# solve equation (2) for A given those four parameters. Output A to full precision.
%# Hint: First think about what dV is for a spherical integral!

\begin{equation}\label{eqn:one}
n(x) = A \langle N_{sat} \rangle \left( \frac{x}{b} \right)^{a-3}
\exp{\left[ -\left(\frac{x}{b}\right)^c \right]}
\end{equation}

For a spherical integral, $\mathrm{d}V = x^2 \sin{\theta} \,\mathrm{d}x \mathrm{d}\theta \mathrm{d}\phi$. When integrating over a full sphere, this simplifies to $4\pi x^2 \,\mathrm{d}x$

\begin{equation}
\begin{split}
\langle N_{sat} \rangle & = \int\int\int_V n(x) \,\mathrm{d}V \\
& = \int_0^{2\pi} \int_0^\pi \int_{x=0}^{x=5} n(x) x^2 \sin{\theta} \,\mathrm{d}x \mathrm{d}\theta \mathrm{d}\phi \\
& = 4\pi \int_{x=0}^{x=5} n(x) x^2 \,\mathrm{d}x
\end{split}
\end{equation}


Filling in \ref{eqn:one},
\begin{equation}
\begin{split}
\frac1A & = 4\pi \int_{x=0}^{x=5} 
\left( \frac{x}{b} \right)^{a-3}
\exp{\left[ -\left(\frac{x}{b}\right)^c \right]}
x^2 \,\mathrm{d}x
\end{split}
\end{equation}

The computed value of $A$ is:
\lstinputlisting{question1a_output.txt}

\subsection{[ding met rng]}
%# We want to generate 3D satellite positions such that they statistically 
%# follow the satellite profile in equation (1); 
%# that is, the probability distribution of the (relative)
%# radii x in [0, 5) should be p(x) dx = N (x) dx/ hN sat i.
%# Use one of the methods discussed in class to sample this distribution.


[[[ transformation kan niet zomaar want N(x) is niet invertible ]]]

$N(x)\mathrm{d}x$ is the number of galaxies in range $[x,x+\mathrm{d}x)$.

The number of galaxies in range $[x,x+\mathrm{d}x)$ is
\begin{equation}
\begin{split}
N(x) \mathrm{d}x & = 4\pi \int_{x}^{x+\mathrm{d}x} x^2 n(x) \,\mathrm{d}x
\end{split}
\end{equation}
so 
\begin{equation}
\begin{split}
N(x) & = 4\pi x^2 n(x)
\end{split}
\end{equation}



\begin{figure}[h!]\label{fig:question1b}
  \centering
  \includegraphics[width=0.9\linewidth]{./plots/question1b.png}
  \caption{
  [[[caption]]]
  }
\end{figure}



\subsection{[selecting galaxies]}
%#Select 100 random satellite galaxies from (b) in a way that is guaranteed to:
%# 1. select every galaxy with equal probability;
%# 2. not draw the same galaxy twice;
%# 3. not reject any draw.
%# Next sort the 100 drawn galaxies from smallest to largest radius
%# plot the number of galaxies within a radius, use a xlog plot from x = 10 -4 to x = x max .

\begin{figure}[h!]\label{fig:question1c}
  \centering
  \includegraphics[width=0.9\linewidth]{./plots/question1c.png}
  \caption{
  [[[caption]]]
  }
\end{figure}

\subsection{[derivative]}
%# Numerically calculate dn(x)/dx at x = 1.
%# Output the value found alongside the analytical result, both to at least 12 significant digits. 
%# Choose your differentiation algorithm such that these are as close as possible.

\begin{equation}
\begin{split}
n(x) & = A \langle N_{sat} \rangle \left( \frac{x}{b} \right)^{a-3}
\exp{\left[ -\left(\frac{x}{b}\right)^c \right]} \\
& = A \langle N_{sat} \rangle 
f(x) g(x)
\end{split}
\end{equation}


\begin{equation}
\begin{split}
\frac{\mathrm{d}n(x)}{\mathrm{d}x} & = \frac{\mathrm{d}}{\mathrm{d}x} A \langle N_{sat} \rangle \left( \frac{x}{b} \right)^{a-3}
\exp{\left[ -\left(\frac{x}{b}\right)^c \right]} \\
& = \frac{\mathrm{d}}{\mathrm{d}x} A \langle N_{sat} \rangle 
f(x) g(x) \\
& = A \langle N_{sat} \rangle 
\left( \frac{\mathrm{d}f(x)}{\mathrm{d}x} g(x) + f(x) \frac{\mathrm{d}g(x)}{\mathrm{d}x} \right)
 \\
\end{split}
\end{equation}

\begin{equation}
\begin{split}
\frac{\mathrm{d}f(x)}{\mathrm{d}x} & = \frac{\mathrm{d}}{\mathrm{d}x}  \left( \frac{x}{b} \right)^{a-3} \\
& = b^{3-a} (a-3) x^{a-4} \\
& = (a-3) \left( \frac{x}{b} \right)^{a-3} x^{-1} \\
& = f(x) (a-3) x^{-1}
\end{split}
\end{equation}

\begin{equation}
\begin{split}
\frac{\mathrm{d}g(x)}{\mathrm{d}x} & = \frac{\mathrm{d}}{\mathrm{d}x} 
\exp{\left[ -\left(\frac{x}{b}\right)^c \right]} \\
& = \exp{\left[ -\left(\frac{x}{b}\right)^c \right]}  \cdot \frac{\mathrm{d}}{\mathrm{d}x} 
\left[ -\left(\frac{x}{b}\right)^c \right] \\
& = \exp{\left[ -\left(\frac{x}{b}\right)^c \right]}  \cdot 
-b^{-c} c x^{c-1} \\
& = \exp{\left[ -\left(\frac{x}{b}\right)^c \right]}  \cdot 
- \left(\frac{x}{b}\right)^c  c x^{-1} \\
& = g(x)  \cdot 
- \left(\frac{x}{b}\right)^c  c x^{-1} \\
\end{split}
\end{equation}

Combine into product rule:
\begin{equation}
\begin{split}
\frac{\mathrm{d}f(x)}{\mathrm{d}x} g(x) + f(x) \frac{\mathrm{d}g(x)}{\mathrm{d}x} 
& = 
f(x) g(x) (a-3) x^{-1} 
+ 
f(x) g(x)  \cdot - \left(\frac{x}{b}\right)^c  c x^{-1} \\
& =
f(x) g(x)
\left[
(a-3) x^{-1}
- \left(\frac{x}{b}\right)^c  c x^{-1}
\right] \\
& =
f(x) g(x)
\left[ a-3 - c \left(\frac{x}{b}\right)^c  \right] x^{-1}
\\
\end{split}
\end{equation}



So
\begin{equation}
\begin{split}
\frac{\mathrm{d}n(x)}{\mathrm{d}x} & = 
A \langle N_{sat} \rangle
f(x) g(x)
\left[ a-3 - c \left(\frac{x}{b}\right)^c  \right] x^{-1} \\
& = n(x)
\left[ a-3 - c \left(\frac{x}{b}\right)^c  \right] x^{-1}
\end{split}
\end{equation}


Output:
\lstinputlisting{question1d_output.txt}


\newpage
All this is implemented in the following script:
\lstinputlisting{question1.py}


